\documentclass[10pt,a4paper]{amsart}
\usepackage{amsmath,amssymb,amsthm,graphicx}
\usepackage{verbatim}
\newenvironment{reference}{\comment}{\endcomment}
%\newenvironment{reference}{}{}
\newenvironment{slogan}{\comment}{\endcomment}
\newenvironment{history}{\comment}{\endcomment}
% For commutative diagrams we use Xy-pic
\usepackage[all]{xy}

% We use 2cell for 2-commutative diagrams.
\xyoption{2cell}
% Theorem environments.
%
\theoremstyle{plain}
\newtheorem{theorem}[subsection]{Theorem}
\newtheorem{proposition}[subsection]{Proposition}
\newtheorem{lemma}[subsection]{Lemma}

\theoremstyle{definition}
\newtheorem{definition}[subsection]{Definition}
\newtheorem{example}[subsection]{Example}
\newtheorem{exercise}[subsection]{Exercise}
\newtheorem{situation}[subsection]{Situation}

\theoremstyle{remark}
\newtheorem{remark}[subsection]{Remark}
\newtheorem{remarks}[subsection]{Remarks}

\numberwithin{equation}{subsection}

% Macros
%
\def\lim{\mathop{\mathrm{lim}}\nolimits}
\def\colim{\mathop{\mathrm{colim}}\nolimits}
\def\Spec{\mathop{\mathrm{Spec}}}
\def\Hom{\mathop{\mathrm{Hom}}\nolimits}
\def\Ext{\mathop{\mathrm{Ext}}\nolimits}
\def\SheafHom{\mathop{\mathcal{H}\!\mathit{om}}\nolimits}
\def\SheafExt{\mathop{\mathcal{E}\!\mathit{xt}}\nolimits}
\def\Sch{\mathit{Sch}}
\def\Mor{\mathop{\mathrm{Mor}}\nolimits}
\def\Ob{\mathop{\mathrm{Ob}}\nolimits}
\def\Sh{\mathop{\mathit{Sh}}\nolimits}
\def\NL{\mathop{N\!L}\nolimits}
\def\CH{\mathop{\mathrm{CH}}\nolimits}
\def\proetale{{pro\text{-}\acute{e}tale}}
\def\etale{{\acute{e}tale}}
\def\QCoh{\mathit{QCoh}}
\def\Ker{\mathop{\mathrm{Ker}}}
\def\Im{\mathop{\mathrm{Im}}}
\def\Coker{\mathop{\mathrm{Coker}}}
\def\Coim{\mathop{\mathrm{Coim}}}

% Boxtimes
%
\DeclareMathSymbol{\boxtimes}{\mathbin}{AMSa}{"02}

%
% Macros for moduli stacks/spaces
%
\def\QCohstack{\mathcal{QC}\!\mathit{oh}}
\def\Cohstack{\mathcal{C}\!\mathit{oh}}
\def\Spacesstack{\mathcal{S}\!\mathit{paces}}
\def\Quotfunctor{\mathrm{Quot}}
\def\Hilbfunctor{\mathrm{Hilb}}
\def\Curvesstack{\mathcal{C}\!\mathit{urves}}
\def\Polarizedstack{\mathcal{P}\!\mathit{olarized}}
\def\Complexesstack{\mathcal{C}\!\mathit{omplexes}}
% \Pic is the operator that assigns to X its picard group, usage \Pic(X)
% \Picardstack_{X/B} denotes the Picard stack of X over B
% \Picardfunctor_{X/B} denotes the Picard functor of X over B
\def\Pic{\mathop{\mathrm{Pic}}\nolimits}
\def\Picardstack{\mathcal{P}\!\mathit{ic}}
\def\Picardfunctor{\mathrm{Pic}}
\def\Deformationcategory{\mathcal{D}\!\mathit{ef}}
\usepackage{cleveref}
\crefname{originaldefinition}{original definition}{original definitions}
\Crefname{originaldefinition}{Original definition}{Original definitions}
\makeatletter
\newlabel{definition-cap-c1}{{02SO}{source}{}{sourcepointer-1}{}}
\newlabel{definition-cap-c1@cref}{{[originaldefinition][0][]02SO}{[1][0][]source}{}{}{}}
\makeatother
\crefname{originaldefinition}{original definition}{original definitions}
\Crefname{originaldefinition}{Original definition}{Original definitions}
\makeatletter
\newlabel{definition-rational-equivalence}{{02RW}{source}{}{sourcepointer-2}{}}
\newlabel{definition-rational-equivalence@cref}{{[originaldefinition][0][]02RW}{[1][0][]source}{}{}{}}
\makeatother
\crefname{originallemma}{original lemma}{original lemmas}
\Crefname{originallemma}{Original lemma}{Original lemmas}
\makeatletter
\newlabel{lemma-c1-cap-additive}{{02SP}{source}{}{sourcepointer-3}{}}
\newlabel{lemma-c1-cap-additive@cref}{{[originallemma][0][]02SP}{[1][0][]source}{}{}{}}
\makeatother
\crefname{originallemma}{original lemma}{original lemmas}
\Crefname{originallemma}{Original lemma}{Original lemmas}
\makeatletter
\newlabel{lemma-commutativity-on-integral}{{02TH}{source}{}{sourcepointer-4}{}}
\newlabel{lemma-commutativity-on-integral@cref}{{[originallemma][0][]02TH}{[1][0][]source}{}{}{}}
\makeatother
\crefname{originallemma}{original lemma}{original lemmas}
\Crefname{originallemma}{Original lemma}{Original lemmas}
\makeatletter
\newlabel{lemma-pushforward-cap-c1}{{02SU}{source}{}{sourcepointer-5}{}}
\newlabel{lemma-pushforward-cap-c1@cref}{{[originallemma][0][]02SU}{[1][0][]source}{}{}{}}
\makeatother
\crefname{originalreference}{original reference}{original references}
\Crefname{originalreference}{Original reference}{Original references}
\makeatletter
\newlabel{situation-setup}{{02QL}{source}{}{sourcepointer-6}{}}
\newlabel{situation-setup@cref}{{[originalreference][0][]02QL}{[1][0][]source}{}{}{}}
\makeatother
\title{The operation of capping with a first Chern class factors through rational equivalence}
\author{The Stacks Project Authors}
\date{stacks-project:chow.tex:local-snapshot-20261012; GNU FDL 1.2 or later}
\begin{document}
\maketitle
\noindent\textit{Editorial note.} All delivered mathematical text is the author's own, copied from the pinned Stacks Project chapter file chow.tex; no statement or proof step was rewritten, reordered or removed.\par\medskip
\noindent\textit{Editorial note.} Preamble provenance: the pinned chapter file opens with \textbackslash{}input\{preamble\} and carries no preamble text of its own, so the project's single Stacks-wide preamble staged at sources/stacks-chow/preamble.tex is excerpted verbatim as three spans, lines 8-12 (the verbatim package together with the project's own reference, slogan and history display-environment definitions, including the author's commented-out older form of the reference environment), lines 14-18 (Xy-pic setup) and lines 154-217 (the native theorem declarations and mathematical macros). The lines 8-12 span is required because a counted statement or its surrounding text may use those display environments, which the project typesets as comments. Nothing was added, removed or redefined; the wrapper class and the amsmath/amssymb/amsthm/graphicx base packages are supplied by the assembler.\par\medskip
\noindent\textit{Editorial note.} Same-chapter pointers are located at their \textbackslash{}label lines inside the pinned chapter file; cross-chapter pointers are located in the project's own tag list staged in the snapshot (sources/stacks-chow/tags.txt), at the verbatim lines carrying their canonical tags.\par\medskip
\noindent\textit{Allowed standard prerequisites (editorial).}\begin{itemize}
\item The definition of capping with a first Chern class and of rational equivalence. Same chapter chow.tex, labels definition-cap-c1 (proof line 4962) and definition-rational-equivalence (proof line 4963). Fixes the operation whose descent is being proved and rewrites the triviality of the cycle as a locally finite sum of principal divisors.
\item The projection formula for capping with a first Chern class. Same chapter chow.tex, label lemma-pushforward-cap-c1, invoked at proof line 4973. Moves the cap through the proper morphism from the disjoint union of the components, reducing the statement to a single integral scheme.
\item Commutativity on an integral scheme and additivity in the invertible sheaf. Same chapter chow.tex, labels lemma-commutativity-on-integral (proof line 4985) and lemma-c1-cap-additive (proof line 4990). The first rewrites the cap with the Chern class of L against the divisor of f as the cap with the trivial invertible sheaf against the divisor of a meromorphic section, and the second shows that this last expression vanishes.
\end{itemize}\medskip
\ifdefined\hypertarget\hypertarget{sourcepointer-1}{}\fi
\noindent\textit{Source reference.} Original definition 02SO: sources/stacks-chow/pinned-original-source.tex, source lines 4284--4284. Navigation pointer for \textbackslash{}ref\{definition-cap-c1\} used inside the counted span (line 4962); the argument that consumes it is retained in the span.\par\medskip
\ifdefined\hypertarget\hypertarget{sourcepointer-2}{}\fi
\noindent\textit{Source reference.} Original definition 02RW: sources/stacks-chow/pinned-original-source.tex, source lines 2882--2882. Navigation pointer for \textbackslash{}ref\{definition-rational-equivalence\} used inside the counted span (line 4963); the argument that consumes it is retained in the span.\par\medskip
\ifdefined\hypertarget\hypertarget{sourcepointer-3}{}\fi
\noindent\textit{Source reference.} Original lemma 02SP: sources/stacks-chow/pinned-original-source.tex, source lines 4330--4330. Navigation pointer for \textbackslash{}ref\{lemma-c1-cap-additive\} used inside the counted span (line 4990); the argument that consumes it is retained in the span.\par\medskip
\ifdefined\hypertarget\hypertarget{sourcepointer-4}{}\fi
\noindent\textit{Source reference.} Original lemma 02TH: sources/stacks-chow/pinned-original-source.tex, source lines 4924--4924. Navigation pointer for \textbackslash{}ref\{lemma-commutativity-on-integral\} used inside the counted span (line 4985); the argument that consumes it is retained in the span.\par\medskip
\ifdefined\hypertarget\hypertarget{sourcepointer-5}{}\fi
\noindent\textit{Source reference.} Original lemma 02SU: sources/stacks-chow/pinned-original-source.tex, source lines 4631--4631. Navigation pointer for \textbackslash{}ref\{lemma-pushforward-cap-c1\} used inside the counted span (line 4973); the argument that consumes it is retained in the span.\par\medskip
\ifdefined\hypertarget\hypertarget{sourcepointer-6}{}\fi
\noindent\textit{Source reference.} Original reference 02QL: sources/stacks-chow/pinned-original-source.tex, source lines 1372--1372. Navigation pointer for \textbackslash{}ref\{situation-setup\} used inside the counted span (line 4949); the argument that consumes it is retained in the span.\par\medskip
\begin{lemma}
\label{lemma-factors}
Let $(S, \delta)$ be as in Situation \ref{situation-setup}.
Let $X$ be locally of finite type over $S$.
Let $\mathcal{L}$ be invertible on $X$.
The operation $\alpha \mapsto c_1(\mathcal{L}) \cap \alpha$
factors through rational equivalence to give an operation
$$
c_1(\mathcal{L}) \cap - : \CH_{k + 1}(X) \to \CH_k(X)
$$
\end{lemma}

\begin{proof}
Let $\alpha \in Z_{k + 1}(X)$, and $\alpha \sim_{rat} 0$.
We have to show that $c_1(\mathcal{L}) \cap \alpha$
as defined in Definition \ref{definition-cap-c1} is zero.
By Definition \ref{definition-rational-equivalence} there
exists a locally finite family $\{W_j\}$ of integral closed
subschemes with $\dim_\delta(W_j) = k + 2$ and rational functions
$f_j \in R(W_j)^*$ such that
$$
\alpha = \sum (i_j)_*\text{div}_{W_j}(f_j)
$$
Note that $p : \coprod W_j \to X$ is a proper morphism,
and hence $\alpha = p_*\alpha'$ where $\alpha' \in Z_{k + 1}(\coprod W_j)$
is the sum of the principal divisors $\text{div}_{W_j}(f_j)$.
By Lemma \ref{lemma-pushforward-cap-c1} we have
$c_1(\mathcal{L}) \cap \alpha = p_*(c_1(p^*\mathcal{L}) \cap \alpha')$.
Hence it suffices to show that each
$c_1(\mathcal{L}|_{W_j}) \cap \text{div}_{W_j}(f_j)$ is zero.
In other words we may assume that $X$ is integral and
$\alpha = \text{div}_X(f)$ for some $f \in R(X)^*$.

\medskip\noindent
Assume $X$ is integral and $\alpha = \text{div}_X(f)$ for some $f \in R(X)^*$.
We can think of $f$ as a regular meromorphic section of the invertible
sheaf $\mathcal{N} = \mathcal{O}_X$. Choose a meromorphic section
$s$ of $\mathcal{L}$ and denote $\beta = \text{div}_\mathcal{L}(s)$.
By Lemma \ref{lemma-commutativity-on-integral}
we conclude that
$$
c_1(\mathcal{L}) \cap \alpha = c_1(\mathcal{O}_X) \cap \beta.
$$
However, by Lemma \ref{lemma-c1-cap-additive} we see that the right hand side
is zero in $\CH_k(X)$ as desired.
\end{proof}

\end{document}
